% Card: the short side (rook endings)
\paragraph{Short side}\index{Rook endings!Short side}

The narrower part is the
\vocab[rook endings]{short-side}{short side}{The narrower of the two parts into
which a pawn (other than a rook's pawn) divides the board. The defending king belongs there,
where it does not stand in the way of its own rook}.
In \xref{diagram}{side-checks} the e-pawn has three files on its right (f, g, h).

The defending king belongs on the short side. There it stays near the pawn without standing in
the way of its own rook, which gives side checks from the \vocab{long-side}{long side}{}.
